\documentclass[11pt,letterpaper]{article}
\usepackage[technical-reference]{cornell-notes}
\usepackage{needspace}

\title{Clause 24: Containers Library - Cornell Notes}
\author{}
\date{}
\setDocTitle{Clause 24: Containers Library}
\setDocSubtitle{C++ 2024 Cornell Notes}
\setDocOwner{C++ 2024 Cornell Notes}
\setDocAuthor{}
\setDocDate{}
\setCornellCollection{C++ 2024 Cornell Notes}
\setCornellUnitType{Clause}
\setCornellUnitNumber{24}
\setCornellUnitTitle{Containers Library}
\setCornellCompanionLabel{C++ Technical Reference}
\hypersetup{pdftitle={Clause 24: Containers Library - Cornell Notes},pdfsubject={C++ technical study notes},pdfauthor={}}

\begin{document}
\maketitle
\begin{CornellReferenceOrientationBox}
\textbf{Purpose.} Defines sequence and associative containers, unordered containers, views and adaptors, container requirements, iterator invalidation, allocator use, complexity, and specialized containers such as arrays, spans, and multidimensional views.
\par\textbf{Role.} The standard model for owning and non-owning collections.
\par\textbf{Principal dependencies.} Uses the library-wide rules of Clause 16 together with the applicable core-language concepts and supporting library clauses.
\par\textbf{Primary audience.} programmers, allocator authors, and container implementers.
\par\textbf{Major subject groups.} General; Requirements; Sequence containers; Associative containers; Unordered associative containers; Container adaptors; Views.
\end{CornellReferenceOrientationBox}
\section{Learning Objectives}
\begin{itemize}[label=\textcolor{CornellReferenceTeal}{$\blacktriangleright$}]
\item Define the governing ideas of General and apply its stated contract at a call site.
\item Identify the governing ideas of Requirements and apply its stated contract at a call site.
\item Distinguish the governing ideas of Sequence containers and apply its stated contract at a call site.
\item Explain the governing ideas of Associative containers and apply its stated contract at a call site.
\item Trace the governing ideas of Unordered associative containers and apply its stated contract at a call site.
\item Apply the governing ideas of Container adaptors and apply its stated contract at a call site.
\item Analyze the governing ideas of Views and apply its stated contract at a call site.
\item Evaluate the governing ideas of Preamble and apply its stated contract at a call site.
\item Diagnose the governing ideas of General containers and apply its stated contract at a call site.
\item Compare the governing ideas of Container data races and apply its stated contract at a call site.
\item Verify the governing ideas of Sequence containers and apply its stated contract at a call site.
\item Relate the governing ideas of Node handles and apply its stated contract at a call site.
\item Classify the governing ideas of Insert return type and apply its stated contract at a call site.
\item Justify the governing ideas of Associative containers and apply its stated contract at a call site.
\item Audit the governing ideas of Unordered associative containers and apply its stated contract at a call site.
\item Summarize the governing ideas of General and apply its stated contract at a call site.
\end{itemize}
\section{Normative Structure Map}
\small\begin{longtable}{@{}>{\RaggedRight\arraybackslash}p{0.13\textwidth}>{\RaggedRight\arraybackslash}p{0.27\textwidth}>{\RaggedRight\arraybackslash}p{0.19\textwidth}>{\RaggedRight\arraybackslash}p{0.31\textwidth}@{}}
\toprule\textbf{Subclause} & \textbf{Subject} & \textbf{Rule category} & \textbf{Key dependencies / entities}\\\midrule\endfirsthead
\toprule\textbf{Subclause} & \textbf{Subject} & \textbf{Rule category} & \textbf{Key dependencies / entities}\\\midrule\endhead
\bottomrule\endfoot
24.1 & General & library semantics & Uses the library-wide rules of Clause 16 together with the applicable core-language concepts and supporting library clauses.\\
24.2 & Requirements & requirements & Preamble, General containers, Container data races, Sequence containers, and related subclauses\\
24.3 & Sequence containers & library semantics & General, Header <array> synopsis, Header <deque> synopsis, Header <forward\_\allowbreak{}list> synopsis, and related subclauses\\
24.4 & Associative containers & library semantics & General, Header <map> synopsis, Header <set> synopsis, Class template map, and related subclauses\\
24.5 & Unordered associative containers & library semantics & General, Header <unordered\_\allowbreak{}map> synopsis, Header <unordered\_\allowbreak{}set> synopsis, Class template unordered\_\allowbreak{}map, and related subclauses\\
24.6 & Container adaptors & library semantics & General, Header <queue> synopsis, Header <stack> synopsis, Header <flat\_\allowbreak{}map> synopsis, and related subclauses\\
24.7 & Views & library semantics & General, Contiguous access, Multidimensional access\\
\end{longtable}\normalsize
\begin{CornellReferenceNormativeBox}A heading maps the clause; it does not replace the detailed conditions. For each operation, read requirements, effects, result, exceptions, complexity, remarks, and notes with their stated normative force.\end{CornellReferenceNormativeBox}
\subsection*{Rule Classification Guide}
\small\begin{longtable}{@{}>{\RaggedRight\arraybackslash}p{0.25\textwidth}>{\RaggedRight\arraybackslash}p{0.67\textwidth}@{}}\toprule\textbf{Label or category} & \textbf{How to read it}\\\midrule\endfirsthead\toprule\textbf{Label or category} & \textbf{How to read it (continued)}\\\midrule\endhead\bottomrule\endfoot
Constraints & Conditions controlling whether a declaration, overload, or specialization participates in the relevant context.\\
Mandates & Requirements checked for the described declaration or specialization; failure has the stated well-formedness consequence.\\
Preconditions & Obligations the caller establishes before evaluation; violating one forfeits the operation's contract.\\
Effects / Returns / Postconditions & Separate the state change, produced result, and state guaranteed after successful completion.\\
Throws / Error conditions & State how failure is reported and which exceptional paths are permitted or required.\\
Complexity & Bounds or exact counts for the operations named by the specification.\\
Remarks / Recommended practice & Remarks refine applicability or participation; recommended practice guides implementations without silently becoming a program requirement.\\
Exposition-only & A device for describing semantics, not a public name on which a program can depend.\\
\end{longtable}\normalsize
\section{Cornell Cue-and-Notes}
\Needspace{8\baselineskip}
\subsection*{24.1 General}
\begin{longtable}{@{}>{\RaggedRight\arraybackslash\columncolor{CornellReferenceCueGray}}p{0.28\textwidth}>{\RaggedRight\arraybackslash}p{0.64\textwidth}@{}}
\rowcolor{CornellReferenceNavy}\color{white}\textbf{Cue / recall prompt} & \color{white}\textbf{Notes}\\\endfirsthead
\rowcolor{CornellReferenceNavy}\color{white}\textbf{Cue / recall prompt} & \color{white}\textbf{Notes (continued)}\\\endhead
\arrayrulecolor{CornellReferenceLineGray}
\textbf{What contract governs General?}\\[-1mm]{\scriptsize\CornellClauseRef{24.1}} & Frames the terminology, applicability, and relationships needed to interpret General; detailed rules remain in the following subclauses.\\\hline
\end{longtable}
\begin{CornellReferenceSummaryBox}General supplies the library semantics portion of this clause. The key review move is to connect its local wording to the clause-wide model, then classify each condition by normative force before relying on it.\end{CornellReferenceSummaryBox}
\Needspace{8\baselineskip}
\subsection*{24.2 Requirements}
\begin{longtable}{@{}>{\RaggedRight\arraybackslash\columncolor{CornellReferenceCueGray}}p{0.28\textwidth}>{\RaggedRight\arraybackslash}p{0.64\textwidth}@{}}
\rowcolor{CornellReferenceNavy}\color{white}\textbf{Cue / recall prompt} & \color{white}\textbf{Notes}\\\endfirsthead
\rowcolor{CornellReferenceNavy}\color{white}\textbf{Cue / recall prompt} & \color{white}\textbf{Notes (continued)}\\\endhead
\arrayrulecolor{CornellReferenceLineGray}
\textbf{What contract governs Preamble?}\\[-1mm]{\scriptsize\CornellClauseRef{24.2.1}} & Frames the terminology, applicability, and relationships needed to interpret Preamble; detailed rules remain in the following subclauses.\\\hline
\textbf{What contract governs General containers?}\\[-1mm]{\scriptsize\CornellClauseRef{24.2.2}} & Specifies the facility family General containers. Read its declarations together with mandates, preconditions, effects, results, exceptions, complexity, and lifetime rules.\\\hline
\textbf{What contract governs Container data races?}\\[-1mm]{\scriptsize\CornellClauseRef{24.2.3}} & Specifies collection behavior for Container data races; check element and allocator requirements, iterator/reference invalidation, exception safety, and complexity per operation.\\\hline
\textbf{What contract governs Sequence containers?}\\[-1mm]{\scriptsize\CornellClauseRef{24.2.4}} & Specifies the facility family Sequence containers. Read its declarations together with mandates, preconditions, effects, results, exceptions, complexity, and lifetime rules.\\\hline
\textbf{What contract governs Node handles?}\\[-1mm]{\scriptsize\CornellClauseRef{24.2.5}} & Specifies the facility family Node handles. Read its declarations together with mandates, preconditions, effects, results, exceptions, complexity, and lifetime rules.\\\hline
\textbf{What contract governs Insert return type?}\\[-1mm]{\scriptsize\CornellClauseRef{24.2.6}} & Specifies the facility family Insert return type. Read its declarations together with mandates, preconditions, effects, results, exceptions, complexity, and lifetime rules.\\\hline
\textbf{What contract governs Associative containers?}\\[-1mm]{\scriptsize\CornellClauseRef{24.2.7}} & Specifies the facility family Associative containers. Read its declarations together with mandates, preconditions, effects, results, exceptions, complexity, and lifetime rules.\\\hline
\textbf{What contract governs Unordered associative containers?}\\[-1mm]{\scriptsize\CornellClauseRef{24.2.8}} & Specifies the facility family Unordered associative containers. Read its declarations together with mandates, preconditions, effects, results, exceptions, complexity, and lifetime rules.\\\hline
\end{longtable}
\begin{CornellReferenceSummaryBox}Requirements supplies the requirements portion of this clause. The key review move is to connect its local wording to the clause-wide model, then classify each condition by normative force before relying on it.\end{CornellReferenceSummaryBox}
\Needspace{5\baselineskip}\paragraph{Detailed coverage ledger.}
\begin{description}[style=nextline,leftmargin=2.8cm,font=\normalfont\bfseries\color{CornellReferenceTeal}]
\item[24.2.2.1] \textbf{Introduction.} Frames the terminology, applicability, and relationships needed to interpret Introduction; detailed rules remain in the following subclauses.
\item[24.2.2.2] \textbf{Container requirements.} States the admissibility and semantic obligations for Container requirements. Separate compile-time participation rules from caller preconditions and runtime effects.
\item[24.2.2.3] \textbf{Reversible container requirements.} States the admissibility and semantic obligations for Reversible container requirements. Separate compile-time participation rules from caller preconditions and runtime effects.
\item[24.2.2.4] \textbf{Optional container requirements.} States the admissibility and semantic obligations for Optional container requirements. Separate compile-time participation rules from caller preconditions and runtime effects.
\item[24.2.2.5] \textbf{Allocator-aware containers.} Specifies the facility family Allocator-aware containers. Read its declarations together with mandates, preconditions, effects, results, exceptions, complexity, and lifetime rules.
\item[24.2.5.1] \textbf{Overview.} Frames the terminology, applicability, and relationships needed to interpret Overview; detailed rules remain in the following subclauses.
\item[24.2.5.2] \textbf{Constructors, copy, and assignment.} Specifies how Constructors, copy, and assignment establishes object state, including participation constraints, initialization effects, and exceptional exits.
\item[24.2.5.3] \textbf{Destructor.} Specifies how Destructor ends the object's managed state and which cleanup or termination consequences apply.
\item[24.2.5.4] \textbf{Observers.} Specifies the facility family Observers. Read its declarations together with mandates, preconditions, effects, results, exceptions, complexity, and lifetime rules.
\item[24.2.5.5] \textbf{Modifiers.} Specifies the facility family Modifiers. Read its declarations together with mandates, preconditions, effects, results, exceptions, complexity, and lifetime rules.
\item[24.2.7.1] \textbf{General.} Frames the terminology, applicability, and relationships needed to interpret General; detailed rules remain in the following subclauses.
\item[24.2.7.2] \textbf{Exception safety guarantees.} Defines the failure model for Exception safety guarantees, including how an error is represented, reported, propagated, or converted into a diagnostic consequence.
\item[24.2.8.1] \textbf{General.} Frames the terminology, applicability, and relationships needed to interpret General; detailed rules remain in the following subclauses.
\item[24.2.8.2] \textbf{Exception safety guarantees.} Defines the failure model for Exception safety guarantees, including how an error is represented, reported, propagated, or converted into a diagnostic consequence.
\end{description}
\Needspace{8\baselineskip}
\subsection*{24.3 Sequence containers}
\begin{longtable}{@{}>{\RaggedRight\arraybackslash\columncolor{CornellReferenceCueGray}}p{0.28\textwidth}>{\RaggedRight\arraybackslash}p{0.64\textwidth}@{}}
\rowcolor{CornellReferenceNavy}\color{white}\textbf{Cue / recall prompt} & \color{white}\textbf{Notes}\\\endfirsthead
\rowcolor{CornellReferenceNavy}\color{white}\textbf{Cue / recall prompt} & \color{white}\textbf{Notes (continued)}\\\endhead
\arrayrulecolor{CornellReferenceLineGray}
\textbf{What contract governs General?}\\[-1mm]{\scriptsize\CornellClauseRef{24.3.1}} & Frames the terminology, applicability, and relationships needed to interpret General; detailed rules remain in the following subclauses.\\\hline
\textbf{What contract governs Header <array> synopsis?}\\[-1mm]{\scriptsize\CornellClauseRef{24.3.2}} & Catalogs the declarations made available by Header <array> synopsis. The synopsis is an interface map; the detailed specifications supply constraints and behavior.\\\hline
\textbf{What contract governs Header <deque> synopsis?}\\[-1mm]{\scriptsize\CornellClauseRef{24.3.3}} & Catalogs the declarations made available by Header <deque> synopsis. The synopsis is an interface map; the detailed specifications supply constraints and behavior.\\\hline
\textbf{What contract governs Header <forward\_\allowbreak{}list> synopsis?}\\[-1mm]{\scriptsize\CornellClauseRef{24.3.4}} & Catalogs the declarations made available by Header <forward\_\allowbreak{}list> synopsis. The synopsis is an interface map; the detailed specifications supply constraints and behavior.\\\hline
\textbf{What contract governs Header <list> synopsis?}\\[-1mm]{\scriptsize\CornellClauseRef{24.3.5}} & Catalogs the declarations made available by Header <list> synopsis. The synopsis is an interface map; the detailed specifications supply constraints and behavior.\\\hline
\textbf{What contract governs Header <vector> synopsis?}\\[-1mm]{\scriptsize\CornellClauseRef{24.3.6}} & Catalogs the declarations made available by Header <vector> synopsis. The synopsis is an interface map; the detailed specifications supply constraints and behavior.\\\hline
\textbf{What contract governs Class template array?}\\[-1mm]{\scriptsize\CornellClauseRef{24.3.7}} & Specifies the facility family Class template array. Read its declarations together with mandates, preconditions, effects, results, exceptions, complexity, and lifetime rules.\\\hline
\textbf{What contract governs Class template deque?}\\[-1mm]{\scriptsize\CornellClauseRef{24.3.8}} & Specifies collection behavior for Class template deque; check element and allocator requirements, iterator/reference invalidation, exception safety, and complexity per operation.\\\hline
\textbf{What contract governs Class template forward\_\allowbreak{}list?}\\[-1mm]{\scriptsize\CornellClauseRef{24.3.9}} & Specifies the facility family Class template forward\_\allowbreak{}list. Read its declarations together with mandates, preconditions, effects, results, exceptions, complexity, and lifetime rules.\\\hline
\textbf{What contract governs Class template list?}\\[-1mm]{\scriptsize\CornellClauseRef{24.3.10}} & Specifies collection behavior for Class template list; check element and allocator requirements, iterator/reference invalidation, exception safety, and complexity per operation.\\\hline
\textbf{What contract governs Class template vector?}\\[-1mm]{\scriptsize\CornellClauseRef{24.3.11}} & Specifies collection behavior for Class template vector; check element and allocator requirements, iterator/reference invalidation, exception safety, and complexity per operation.\\\hline
\textbf{What contract governs Specialization of vector for bool?}\\[-1mm]{\scriptsize\CornellClauseRef{24.3.12}} & Specifies collection behavior for Specialization of vector for bool; check element and allocator requirements, iterator/reference invalidation, exception safety, and complexity per operation.\\\hline
\end{longtable}
\begin{CornellReferenceSummaryBox}Sequence containers supplies the library semantics portion of this clause. The key review move is to connect its local wording to the clause-wide model, then classify each condition by normative force before relying on it.\end{CornellReferenceSummaryBox}
\Needspace{5\baselineskip}\paragraph{Detailed coverage ledger.}
\begin{description}[style=nextline,leftmargin=2.8cm,font=\normalfont\bfseries\color{CornellReferenceTeal}]
\item[24.3.7.1] \textbf{Overview.} Frames the terminology, applicability, and relationships needed to interpret Overview; detailed rules remain in the following subclauses.
\item[24.3.7.2] \textbf{Constructors, copy, and assignment.} Specifies how Constructors, copy, and assignment establishes object state, including participation constraints, initialization effects, and exceptional exits.
\item[24.3.7.3] \textbf{Member functions.} Specifies the facility family Member functions. Read its declarations together with mandates, preconditions, effects, results, exceptions, complexity, and lifetime rules.
\item[24.3.7.4] \textbf{Specialized algorithms.} Defines the observable result of Specialized algorithms together with applicable iterator-domain, ordering, complexity, invalidation, and exception conditions.
\item[24.3.7.5] \textbf{Zero-sized arrays.} Specifies the facility family Zero-sized arrays. Read its declarations together with mandates, preconditions, effects, results, exceptions, complexity, and lifetime rules.
\item[24.3.7.6] \textbf{Array creation functions.} Specifies the facility family Array creation functions. Read its declarations together with mandates, preconditions, effects, results, exceptions, complexity, and lifetime rules.
\item[24.3.7.7] \textbf{Tuple interface.} Specifies the facility family Tuple interface. Read its declarations together with mandates, preconditions, effects, results, exceptions, complexity, and lifetime rules.
\item[24.3.8.1] \textbf{Overview.} Frames the terminology, applicability, and relationships needed to interpret Overview; detailed rules remain in the following subclauses.
\item[24.3.8.2] \textbf{Constructors, copy, and assignment.} Specifies how Constructors, copy, and assignment establishes object state, including participation constraints, initialization effects, and exceptional exits.
\item[24.3.8.3] \textbf{Capacity.} Specifies the facility family Capacity. Read its declarations together with mandates, preconditions, effects, results, exceptions, complexity, and lifetime rules.
\item[24.3.8.4] \textbf{Modifiers.} Specifies the facility family Modifiers. Read its declarations together with mandates, preconditions, effects, results, exceptions, complexity, and lifetime rules.
\item[24.3.8.5] \textbf{Erasure.} Specifies the facility family Erasure. Read its declarations together with mandates, preconditions, effects, results, exceptions, complexity, and lifetime rules.
\item[24.3.9.1] \textbf{Overview.} Frames the terminology, applicability, and relationships needed to interpret Overview; detailed rules remain in the following subclauses.
\item[24.3.9.2] \textbf{Constructors, copy, and assignment.} Specifies how Constructors, copy, and assignment establishes object state, including participation constraints, initialization effects, and exceptional exits.
\item[24.3.9.3] \textbf{Iterators.} Explains the traversal or view contract for Iterators; prove domain validity, lifetime, sentinel compatibility, and any borrowing or invalidation property.
\item[24.3.9.4] \textbf{Element access.} Specifies the facility family Element access. Read its declarations together with mandates, preconditions, effects, results, exceptions, complexity, and lifetime rules.
\item[24.3.9.5] \textbf{Modifiers.} Specifies the facility family Modifiers. Read its declarations together with mandates, preconditions, effects, results, exceptions, complexity, and lifetime rules.
\item[24.3.9.6] \textbf{Operations.} Defines the observable result of Operations together with applicable iterator-domain, ordering, complexity, invalidation, and exception conditions.
\item[24.3.9.7] \textbf{Erasure.} Specifies the facility family Erasure. Read its declarations together with mandates, preconditions, effects, results, exceptions, complexity, and lifetime rules.
\item[24.3.10.1] \textbf{Overview.} Frames the terminology, applicability, and relationships needed to interpret Overview; detailed rules remain in the following subclauses.
\item[24.3.10.2] \textbf{Constructors, copy, and assignment.} Specifies how Constructors, copy, and assignment establishes object state, including participation constraints, initialization effects, and exceptional exits.
\item[24.3.10.3] \textbf{Capacity.} Specifies the facility family Capacity. Read its declarations together with mandates, preconditions, effects, results, exceptions, complexity, and lifetime rules.
\item[24.3.10.4] \textbf{Modifiers.} Specifies the facility family Modifiers. Read its declarations together with mandates, preconditions, effects, results, exceptions, complexity, and lifetime rules.
\item[24.3.10.5] \textbf{Operations.} Defines the observable result of Operations together with applicable iterator-domain, ordering, complexity, invalidation, and exception conditions.
\item[24.3.10.6] \textbf{Erasure.} Specifies the facility family Erasure. Read its declarations together with mandates, preconditions, effects, results, exceptions, complexity, and lifetime rules.
\item[24.3.11.1] \textbf{Overview.} Frames the terminology, applicability, and relationships needed to interpret Overview; detailed rules remain in the following subclauses.
\item[24.3.11.2] \textbf{Constructors.} Specifies how Constructors establishes object state, including participation constraints, initialization effects, and exceptional exits.
\item[24.3.11.3] \textbf{Capacity.} Specifies the facility family Capacity. Read its declarations together with mandates, preconditions, effects, results, exceptions, complexity, and lifetime rules.
\item[24.3.11.4] \textbf{Data.} Specifies the facility family Data. Read its declarations together with mandates, preconditions, effects, results, exceptions, complexity, and lifetime rules.
\item[24.3.11.5] \textbf{Modifiers.} Specifies the facility family Modifiers. Read its declarations together with mandates, preconditions, effects, results, exceptions, complexity, and lifetime rules.
\item[24.3.11.6] \textbf{Erasure.} Specifies the facility family Erasure. Read its declarations together with mandates, preconditions, effects, results, exceptions, complexity, and lifetime rules.
\item[24.3.12.1] \textbf{Partial class template specialization vector<bool, Allocator>.} Specifies collection behavior for Partial class template specialization vector<bool, Allocator>; check element and allocator requirements, iterator/reference invalidation, exception safety, and complexity per operation.
\item[24.3.12.2] \textbf{Formatter specialization for vector<bool>.} Specifies collection behavior for Formatter specialization for vector<bool>; check element and allocator requirements, iterator/reference invalidation, exception safety, and complexity per operation.
\end{description}
\Needspace{8\baselineskip}
\subsection*{24.4 Associative containers}
\begin{longtable}{@{}>{\RaggedRight\arraybackslash\columncolor{CornellReferenceCueGray}}p{0.28\textwidth}>{\RaggedRight\arraybackslash}p{0.64\textwidth}@{}}
\rowcolor{CornellReferenceNavy}\color{white}\textbf{Cue / recall prompt} & \color{white}\textbf{Notes}\\\endfirsthead
\rowcolor{CornellReferenceNavy}\color{white}\textbf{Cue / recall prompt} & \color{white}\textbf{Notes (continued)}\\\endhead
\arrayrulecolor{CornellReferenceLineGray}
\textbf{What contract governs General?}\\[-1mm]{\scriptsize\CornellClauseRef{24.4.1}} & Frames the terminology, applicability, and relationships needed to interpret General; detailed rules remain in the following subclauses.\\\hline
\textbf{What contract governs Header <map> synopsis?}\\[-1mm]{\scriptsize\CornellClauseRef{24.4.2}} & Catalogs the declarations made available by Header <map> synopsis. The synopsis is an interface map; the detailed specifications supply constraints and behavior.\\\hline
\textbf{What contract governs Header <set> synopsis?}\\[-1mm]{\scriptsize\CornellClauseRef{24.4.3}} & Catalogs the declarations made available by Header <set> synopsis. The synopsis is an interface map; the detailed specifications supply constraints and behavior.\\\hline
\textbf{What contract governs Class template map?}\\[-1mm]{\scriptsize\CornellClauseRef{24.4.4}} & Specifies collection behavior for Class template map; check element and allocator requirements, iterator/reference invalidation, exception safety, and complexity per operation.\\\hline
\textbf{What contract governs Class template multimap?}\\[-1mm]{\scriptsize\CornellClauseRef{24.4.5}} & Specifies the facility family Class template multimap. Read its declarations together with mandates, preconditions, effects, results, exceptions, complexity, and lifetime rules.\\\hline
\textbf{What contract governs Class template set?}\\[-1mm]{\scriptsize\CornellClauseRef{24.4.6}} & Specifies collection behavior for Class template set; check element and allocator requirements, iterator/reference invalidation, exception safety, and complexity per operation.\\\hline
\textbf{What contract governs Class template multiset?}\\[-1mm]{\scriptsize\CornellClauseRef{24.4.7}} & Specifies the facility family Class template multiset. Read its declarations together with mandates, preconditions, effects, results, exceptions, complexity, and lifetime rules.\\\hline
\end{longtable}
\begin{CornellReferenceSummaryBox}Associative containers supplies the library semantics portion of this clause. The key review move is to connect its local wording to the clause-wide model, then classify each condition by normative force before relying on it.\end{CornellReferenceSummaryBox}
\Needspace{5\baselineskip}\paragraph{Detailed coverage ledger.}
\begin{description}[style=nextline,leftmargin=2.8cm,font=\normalfont\bfseries\color{CornellReferenceTeal}]
\item[24.4.4.1] \textbf{Overview.} Frames the terminology, applicability, and relationships needed to interpret Overview; detailed rules remain in the following subclauses.
\item[24.4.4.2] \textbf{Constructors, copy, and assignment.} Specifies how Constructors, copy, and assignment establishes object state, including participation constraints, initialization effects, and exceptional exits.
\item[24.4.4.3] \textbf{Element access.} Specifies the facility family Element access. Read its declarations together with mandates, preconditions, effects, results, exceptions, complexity, and lifetime rules.
\item[24.4.4.4] \textbf{Modifiers.} Specifies the facility family Modifiers. Read its declarations together with mandates, preconditions, effects, results, exceptions, complexity, and lifetime rules.
\item[24.4.4.5] \textbf{Erasure.} Specifies the facility family Erasure. Read its declarations together with mandates, preconditions, effects, results, exceptions, complexity, and lifetime rules.
\item[24.4.5.1] \textbf{Overview.} Frames the terminology, applicability, and relationships needed to interpret Overview; detailed rules remain in the following subclauses.
\item[24.4.5.2] \textbf{Constructors.} Specifies how Constructors establishes object state, including participation constraints, initialization effects, and exceptional exits.
\item[24.4.5.3] \textbf{Modifiers.} Specifies the facility family Modifiers. Read its declarations together with mandates, preconditions, effects, results, exceptions, complexity, and lifetime rules.
\item[24.4.5.4] \textbf{Erasure.} Specifies the facility family Erasure. Read its declarations together with mandates, preconditions, effects, results, exceptions, complexity, and lifetime rules.
\item[24.4.6.1] \textbf{Overview.} Frames the terminology, applicability, and relationships needed to interpret Overview; detailed rules remain in the following subclauses.
\item[24.4.6.2] \textbf{Constructors, copy, and assignment.} Specifies how Constructors, copy, and assignment establishes object state, including participation constraints, initialization effects, and exceptional exits.
\item[24.4.6.3] \textbf{Erasure.} Specifies the facility family Erasure. Read its declarations together with mandates, preconditions, effects, results, exceptions, complexity, and lifetime rules.
\item[24.4.7.1] \textbf{Overview.} Frames the terminology, applicability, and relationships needed to interpret Overview; detailed rules remain in the following subclauses.
\item[24.4.7.2] \textbf{Constructors.} Specifies how Constructors establishes object state, including participation constraints, initialization effects, and exceptional exits.
\item[24.4.7.3] \textbf{Erasure.} Specifies the facility family Erasure. Read its declarations together with mandates, preconditions, effects, results, exceptions, complexity, and lifetime rules.
\end{description}
\Needspace{8\baselineskip}
\subsection*{24.5 Unordered associative containers}
\begin{longtable}{@{}>{\RaggedRight\arraybackslash\columncolor{CornellReferenceCueGray}}p{0.28\textwidth}>{\RaggedRight\arraybackslash}p{0.64\textwidth}@{}}
\rowcolor{CornellReferenceNavy}\color{white}\textbf{Cue / recall prompt} & \color{white}\textbf{Notes}\\\endfirsthead
\rowcolor{CornellReferenceNavy}\color{white}\textbf{Cue / recall prompt} & \color{white}\textbf{Notes (continued)}\\\endhead
\arrayrulecolor{CornellReferenceLineGray}
\textbf{What contract governs General?}\\[-1mm]{\scriptsize\CornellClauseRef{24.5.1}} & Frames the terminology, applicability, and relationships needed to interpret General; detailed rules remain in the following subclauses.\\\hline
\textbf{What contract governs Header <unordered\_\allowbreak{}map> synopsis?}\\[-1mm]{\scriptsize\CornellClauseRef{24.5.2}} & Catalogs the declarations made available by Header <unordered\_\allowbreak{}map> synopsis. The synopsis is an interface map; the detailed specifications supply constraints and behavior.\\\hline
\textbf{What contract governs Header <unordered\_\allowbreak{}set> synopsis?}\\[-1mm]{\scriptsize\CornellClauseRef{24.5.3}} & Catalogs the declarations made available by Header <unordered\_\allowbreak{}set> synopsis. The synopsis is an interface map; the detailed specifications supply constraints and behavior.\\\hline
\textbf{What contract governs Class template unordered\_\allowbreak{}map?}\\[-1mm]{\scriptsize\CornellClauseRef{24.5.4}} & Specifies the facility family Class template unordered\_\allowbreak{}map. Read its declarations together with mandates, preconditions, effects, results, exceptions, complexity, and lifetime rules.\\\hline
\textbf{What contract governs Class template unordered\_\allowbreak{}multimap?}\\[-1mm]{\scriptsize\CornellClauseRef{24.5.5}} & Specifies the facility family Class template unordered\_\allowbreak{}multimap. Read its declarations together with mandates, preconditions, effects, results, exceptions, complexity, and lifetime rules.\\\hline
\textbf{What contract governs Class template unordered\_\allowbreak{}set?}\\[-1mm]{\scriptsize\CornellClauseRef{24.5.6}} & Specifies the facility family Class template unordered\_\allowbreak{}set. Read its declarations together with mandates, preconditions, effects, results, exceptions, complexity, and lifetime rules.\\\hline
\textbf{What contract governs Class template unordered\_\allowbreak{}multiset?}\\[-1mm]{\scriptsize\CornellClauseRef{24.5.7}} & Specifies the facility family Class template unordered\_\allowbreak{}multiset. Read its declarations together with mandates, preconditions, effects, results, exceptions, complexity, and lifetime rules.\\\hline
\end{longtable}
\begin{CornellReferenceSummaryBox}Unordered associative containers supplies the library semantics portion of this clause. The key review move is to connect its local wording to the clause-wide model, then classify each condition by normative force before relying on it.\end{CornellReferenceSummaryBox}
\Needspace{5\baselineskip}\paragraph{Detailed coverage ledger.}
\begin{description}[style=nextline,leftmargin=2.8cm,font=\normalfont\bfseries\color{CornellReferenceTeal}]
\item[24.5.4.1] \textbf{Overview.} Frames the terminology, applicability, and relationships needed to interpret Overview; detailed rules remain in the following subclauses.
\item[24.5.4.2] \textbf{Constructors.} Specifies how Constructors establishes object state, including participation constraints, initialization effects, and exceptional exits.
\item[24.5.4.3] \textbf{Element access.} Specifies the facility family Element access. Read its declarations together with mandates, preconditions, effects, results, exceptions, complexity, and lifetime rules.
\item[24.5.4.4] \textbf{Modifiers.} Specifies the facility family Modifiers. Read its declarations together with mandates, preconditions, effects, results, exceptions, complexity, and lifetime rules.
\item[24.5.4.5] \textbf{Erasure.} Specifies the facility family Erasure. Read its declarations together with mandates, preconditions, effects, results, exceptions, complexity, and lifetime rules.
\item[24.5.5.1] \textbf{Overview.} Frames the terminology, applicability, and relationships needed to interpret Overview; detailed rules remain in the following subclauses.
\item[24.5.5.2] \textbf{Constructors.} Specifies how Constructors establishes object state, including participation constraints, initialization effects, and exceptional exits.
\item[24.5.5.3] \textbf{Modifiers.} Specifies the facility family Modifiers. Read its declarations together with mandates, preconditions, effects, results, exceptions, complexity, and lifetime rules.
\item[24.5.5.4] \textbf{Erasure.} Specifies the facility family Erasure. Read its declarations together with mandates, preconditions, effects, results, exceptions, complexity, and lifetime rules.
\item[24.5.6.1] \textbf{Overview.} Frames the terminology, applicability, and relationships needed to interpret Overview; detailed rules remain in the following subclauses.
\item[24.5.6.2] \textbf{Constructors.} Specifies how Constructors establishes object state, including participation constraints, initialization effects, and exceptional exits.
\item[24.5.6.3] \textbf{Erasure.} Specifies the facility family Erasure. Read its declarations together with mandates, preconditions, effects, results, exceptions, complexity, and lifetime rules.
\item[24.5.7.1] \textbf{Overview.} Frames the terminology, applicability, and relationships needed to interpret Overview; detailed rules remain in the following subclauses.
\item[24.5.7.2] \textbf{Constructors.} Specifies how Constructors establishes object state, including participation constraints, initialization effects, and exceptional exits.
\item[24.5.7.3] \textbf{Erasure.} Specifies the facility family Erasure. Read its declarations together with mandates, preconditions, effects, results, exceptions, complexity, and lifetime rules.
\end{description}
\Needspace{8\baselineskip}
\subsection*{24.6 Container adaptors}
\begin{longtable}{@{}>{\RaggedRight\arraybackslash\columncolor{CornellReferenceCueGray}}p{0.28\textwidth}>{\RaggedRight\arraybackslash}p{0.64\textwidth}@{}}
\rowcolor{CornellReferenceNavy}\color{white}\textbf{Cue / recall prompt} & \color{white}\textbf{Notes}\\\endfirsthead
\rowcolor{CornellReferenceNavy}\color{white}\textbf{Cue / recall prompt} & \color{white}\textbf{Notes (continued)}\\\endhead
\arrayrulecolor{CornellReferenceLineGray}
\textbf{What contract governs General?}\\[-1mm]{\scriptsize\CornellClauseRef{24.6.1}} & Frames the terminology, applicability, and relationships needed to interpret General; detailed rules remain in the following subclauses.\\\hline
\textbf{What contract governs Header <queue> synopsis?}\\[-1mm]{\scriptsize\CornellClauseRef{24.6.2}} & Catalogs the declarations made available by Header <queue> synopsis. The synopsis is an interface map; the detailed specifications supply constraints and behavior.\\\hline
\textbf{What contract governs Header <stack> synopsis?}\\[-1mm]{\scriptsize\CornellClauseRef{24.6.3}} & Catalogs the declarations made available by Header <stack> synopsis. The synopsis is an interface map; the detailed specifications supply constraints and behavior.\\\hline
\textbf{What contract governs Header <flat\_\allowbreak{}map> synopsis?}\\[-1mm]{\scriptsize\CornellClauseRef{24.6.4}} & Catalogs the declarations made available by Header <flat\_\allowbreak{}map> synopsis. The synopsis is an interface map; the detailed specifications supply constraints and behavior.\\\hline
\textbf{What contract governs Header <flat\_\allowbreak{}set> synopsis?}\\[-1mm]{\scriptsize\CornellClauseRef{24.6.5}} & Catalogs the declarations made available by Header <flat\_\allowbreak{}set> synopsis. The synopsis is an interface map; the detailed specifications supply constraints and behavior.\\\hline
\textbf{What contract governs Class template queue?}\\[-1mm]{\scriptsize\CornellClauseRef{24.6.6}} & Specifies the facility family Class template queue. Read its declarations together with mandates, preconditions, effects, results, exceptions, complexity, and lifetime rules.\\\hline
\textbf{What contract governs Class template priority\_\allowbreak{}queue?}\\[-1mm]{\scriptsize\CornellClauseRef{24.6.7}} & Specifies the facility family Class template priority\_\allowbreak{}queue. Read its declarations together with mandates, preconditions, effects, results, exceptions, complexity, and lifetime rules.\\\hline
\textbf{What contract governs Class template stack?}\\[-1mm]{\scriptsize\CornellClauseRef{24.6.8}} & Specifies the facility family Class template stack. Read its declarations together with mandates, preconditions, effects, results, exceptions, complexity, and lifetime rules.\\\hline
\textbf{What contract governs Class template flat\_\allowbreak{}map?}\\[-1mm]{\scriptsize\CornellClauseRef{24.6.9}} & Specifies the facility family Class template flat\_\allowbreak{}map. Read its declarations together with mandates, preconditions, effects, results, exceptions, complexity, and lifetime rules.\\\hline
\textbf{What contract governs Class template flat\_\allowbreak{}multimap?}\\[-1mm]{\scriptsize\CornellClauseRef{24.6.10}} & Specifies the facility family Class template flat\_\allowbreak{}multimap. Read its declarations together with mandates, preconditions, effects, results, exceptions, complexity, and lifetime rules.\\\hline
\textbf{What contract governs Class template flat\_\allowbreak{}set?}\\[-1mm]{\scriptsize\CornellClauseRef{24.6.11}} & Specifies the facility family Class template flat\_\allowbreak{}set. Read its declarations together with mandates, preconditions, effects, results, exceptions, complexity, and lifetime rules.\\\hline
\textbf{What contract governs Class template flat\_\allowbreak{}multiset?}\\[-1mm]{\scriptsize\CornellClauseRef{24.6.12}} & Specifies the facility family Class template flat\_\allowbreak{}multiset. Read its declarations together with mandates, preconditions, effects, results, exceptions, complexity, and lifetime rules.\\\hline
\textbf{What contract governs Container adaptors formatting?}\\[-1mm]{\scriptsize\CornellClauseRef{24.6.13}} & Specifies collection behavior for Container adaptors formatting; check element and allocator requirements, iterator/reference invalidation, exception safety, and complexity per operation.\\\hline
\end{longtable}
\begin{CornellReferenceSummaryBox}Container adaptors supplies the library semantics portion of this clause. The key review move is to connect its local wording to the clause-wide model, then classify each condition by normative force before relying on it.\end{CornellReferenceSummaryBox}
\Needspace{5\baselineskip}\paragraph{Detailed coverage ledger.}
\begin{description}[style=nextline,leftmargin=2.8cm,font=\normalfont\bfseries\color{CornellReferenceTeal}]
\item[24.6.6.1] \textbf{Definition.} Specifies the facility family Definition. Read its declarations together with mandates, preconditions, effects, results, exceptions, complexity, and lifetime rules.
\item[24.6.6.2] \textbf{Constructors.} Specifies how Constructors establishes object state, including participation constraints, initialization effects, and exceptional exits.
\item[24.6.6.3] \textbf{Constructors with allocators.} Specifies how Constructors with allocators establishes object state, including participation constraints, initialization effects, and exceptional exits.
\item[24.6.6.4] \textbf{Modifiers.} Specifies the facility family Modifiers. Read its declarations together with mandates, preconditions, effects, results, exceptions, complexity, and lifetime rules.
\item[24.6.6.5] \textbf{Operators.} Specifies the facility family Operators. Read its declarations together with mandates, preconditions, effects, results, exceptions, complexity, and lifetime rules.
\item[24.6.6.6] \textbf{Specialized algorithms.} Defines the observable result of Specialized algorithms together with applicable iterator-domain, ordering, complexity, invalidation, and exception conditions.
\item[24.6.7.1] \textbf{Overview.} Frames the terminology, applicability, and relationships needed to interpret Overview; detailed rules remain in the following subclauses.
\item[24.6.7.2] \textbf{Constructors.} Specifies how Constructors establishes object state, including participation constraints, initialization effects, and exceptional exits.
\item[24.6.7.3] \textbf{Constructors with allocators.} Specifies how Constructors with allocators establishes object state, including participation constraints, initialization effects, and exceptional exits.
\item[24.6.7.4] \textbf{Members.} Specifies the facility family Members. Read its declarations together with mandates, preconditions, effects, results, exceptions, complexity, and lifetime rules.
\item[24.6.7.5] \textbf{Specialized algorithms.} Defines the observable result of Specialized algorithms together with applicable iterator-domain, ordering, complexity, invalidation, and exception conditions.
\item[24.6.8.1] \textbf{General.} Frames the terminology, applicability, and relationships needed to interpret General; detailed rules remain in the following subclauses.
\item[24.6.8.2] \textbf{Definition.} Specifies the facility family Definition. Read its declarations together with mandates, preconditions, effects, results, exceptions, complexity, and lifetime rules.
\item[24.6.8.3] \textbf{Constructors.} Specifies how Constructors establishes object state, including participation constraints, initialization effects, and exceptional exits.
\item[24.6.8.4] \textbf{Constructors with allocators.} Specifies how Constructors with allocators establishes object state, including participation constraints, initialization effects, and exceptional exits.
\item[24.6.8.5] \textbf{Modifiers.} Specifies the facility family Modifiers. Read its declarations together with mandates, preconditions, effects, results, exceptions, complexity, and lifetime rules.
\item[24.6.8.6] \textbf{Operators.} Specifies the facility family Operators. Read its declarations together with mandates, preconditions, effects, results, exceptions, complexity, and lifetime rules.
\item[24.6.8.7] \textbf{Specialized algorithms.} Defines the observable result of Specialized algorithms together with applicable iterator-domain, ordering, complexity, invalidation, and exception conditions.
\item[24.6.9.1] \textbf{Overview.} Frames the terminology, applicability, and relationships needed to interpret Overview; detailed rules remain in the following subclauses.
\item[24.6.9.2] \textbf{Definition.} Specifies the facility family Definition. Read its declarations together with mandates, preconditions, effects, results, exceptions, complexity, and lifetime rules.
\item[24.6.9.3] \textbf{Constructors.} Specifies how Constructors establishes object state, including participation constraints, initialization effects, and exceptional exits.
\item[24.6.9.4] \textbf{Capacity.} Specifies the facility family Capacity. Read its declarations together with mandates, preconditions, effects, results, exceptions, complexity, and lifetime rules.
\item[24.6.9.5] \textbf{Access.} Specifies the facility family Access. Read its declarations together with mandates, preconditions, effects, results, exceptions, complexity, and lifetime rules.
\item[24.6.9.6] \textbf{Modifiers.} Specifies the facility family Modifiers. Read its declarations together with mandates, preconditions, effects, results, exceptions, complexity, and lifetime rules.
\item[24.6.9.7] \textbf{Erasure.} Specifies the facility family Erasure. Read its declarations together with mandates, preconditions, effects, results, exceptions, complexity, and lifetime rules.
\item[24.6.10.1] \textbf{Overview.} Frames the terminology, applicability, and relationships needed to interpret Overview; detailed rules remain in the following subclauses.
\item[24.6.10.2] \textbf{Definition.} Specifies the facility family Definition. Read its declarations together with mandates, preconditions, effects, results, exceptions, complexity, and lifetime rules.
\item[24.6.10.3] \textbf{Constructors.} Specifies how Constructors establishes object state, including participation constraints, initialization effects, and exceptional exits.
\item[24.6.10.4] \textbf{Erasure.} Specifies the facility family Erasure. Read its declarations together with mandates, preconditions, effects, results, exceptions, complexity, and lifetime rules.
\item[24.6.11.1] \textbf{Overview.} Frames the terminology, applicability, and relationships needed to interpret Overview; detailed rules remain in the following subclauses.
\item[24.6.11.2] \textbf{Definition.} Specifies the facility family Definition. Read its declarations together with mandates, preconditions, effects, results, exceptions, complexity, and lifetime rules.
\item[24.6.11.3] \textbf{Constructors.} Specifies how Constructors establishes object state, including participation constraints, initialization effects, and exceptional exits.
\item[24.6.11.4] \textbf{Modifiers.} Specifies the facility family Modifiers. Read its declarations together with mandates, preconditions, effects, results, exceptions, complexity, and lifetime rules.
\item[24.6.11.5] \textbf{Erasure.} Specifies the facility family Erasure. Read its declarations together with mandates, preconditions, effects, results, exceptions, complexity, and lifetime rules.
\item[24.6.12.1] \textbf{Overview.} Frames the terminology, applicability, and relationships needed to interpret Overview; detailed rules remain in the following subclauses.
\item[24.6.12.2] \textbf{Definition.} Specifies the facility family Definition. Read its declarations together with mandates, preconditions, effects, results, exceptions, complexity, and lifetime rules.
\item[24.6.12.3] \textbf{Constructors.} Specifies how Constructors establishes object state, including participation constraints, initialization effects, and exceptional exits.
\item[24.6.12.4] \textbf{Modifiers.} Specifies the facility family Modifiers. Read its declarations together with mandates, preconditions, effects, results, exceptions, complexity, and lifetime rules.
\item[24.6.12.5] \textbf{Erasure.} Specifies the facility family Erasure. Read its declarations together with mandates, preconditions, effects, results, exceptions, complexity, and lifetime rules.
\end{description}
\Needspace{8\baselineskip}
\subsection*{24.7 Views}
\begin{longtable}{@{}>{\RaggedRight\arraybackslash\columncolor{CornellReferenceCueGray}}p{0.28\textwidth}>{\RaggedRight\arraybackslash}p{0.64\textwidth}@{}}
\rowcolor{CornellReferenceNavy}\color{white}\textbf{Cue / recall prompt} & \color{white}\textbf{Notes}\\\endfirsthead
\rowcolor{CornellReferenceNavy}\color{white}\textbf{Cue / recall prompt} & \color{white}\textbf{Notes (continued)}\\\endhead
\arrayrulecolor{CornellReferenceLineGray}
\textbf{What contract governs General?}\\[-1mm]{\scriptsize\CornellClauseRef{24.7.1}} & Frames the terminology, applicability, and relationships needed to interpret General; detailed rules remain in the following subclauses.\\\hline
\textbf{What contract governs Contiguous access?}\\[-1mm]{\scriptsize\CornellClauseRef{24.7.2}} & Specifies the facility family Contiguous access. Read its declarations together with mandates, preconditions, effects, results, exceptions, complexity, and lifetime rules.\\\hline
\textbf{What contract governs Multidimensional access?}\\[-1mm]{\scriptsize\CornellClauseRef{24.7.3}} & Specifies the facility family Multidimensional access. Read its declarations together with mandates, preconditions, effects, results, exceptions, complexity, and lifetime rules.\\\hline
\end{longtable}
\begin{CornellReferenceSummaryBox}Views supplies the library semantics portion of this clause. The key review move is to connect its local wording to the clause-wide model, then classify each condition by normative force before relying on it.\end{CornellReferenceSummaryBox}
\Needspace{5\baselineskip}\paragraph{Detailed coverage ledger.}
\begin{description}[style=nextline,leftmargin=2.8cm,font=\normalfont\bfseries\color{CornellReferenceTeal}]
\item[24.7.2.1] \textbf{Header <span> synopsis.} Catalogs the declarations made available by Header <span> synopsis. The synopsis is an interface map; the detailed specifications supply constraints and behavior.
\item[24.7.2.2] \textbf{Class template span.} Specifies the facility family Class template span. Read its declarations together with mandates, preconditions, effects, results, exceptions, complexity, and lifetime rules.
\item[24.7.2.2.1] \textbf{Overview.} Frames the terminology, applicability, and relationships needed to interpret Overview; detailed rules remain in the following subclauses.
\item[24.7.2.2.2] \textbf{Constructors, copy, and assignment.} Specifies how Constructors, copy, and assignment establishes object state, including participation constraints, initialization effects, and exceptional exits.
\item[24.7.2.2.3] \textbf{Deduction guides.} Specifies the facility family Deduction guides. Read its declarations together with mandates, preconditions, effects, results, exceptions, complexity, and lifetime rules.
\item[24.7.2.2.4] \textbf{Subviews.} Explains the traversal or view contract for Subviews; prove domain validity, lifetime, sentinel compatibility, and any borrowing or invalidation property.
\item[24.7.2.2.5] \textbf{Observers.} Specifies the facility family Observers. Read its declarations together with mandates, preconditions, effects, results, exceptions, complexity, and lifetime rules.
\item[24.7.2.2.6] \textbf{Element access.} Specifies the facility family Element access. Read its declarations together with mandates, preconditions, effects, results, exceptions, complexity, and lifetime rules.
\item[24.7.2.2.7] \textbf{Iterator support.} Explains the traversal or view contract for Iterator support; prove domain validity, lifetime, sentinel compatibility, and any borrowing or invalidation property.
\item[24.7.2.3] \textbf{Views of object representation.} Explains the traversal or view contract for Views of object representation; prove domain validity, lifetime, sentinel compatibility, and any borrowing or invalidation property.
\item[24.7.3.1] \textbf{Overview.} Frames the terminology, applicability, and relationships needed to interpret Overview; detailed rules remain in the following subclauses.
\item[24.7.3.2] \textbf{Header <mdspan> synopsis.} Catalogs the declarations made available by Header <mdspan> synopsis. The synopsis is an interface map; the detailed specifications supply constraints and behavior.
\item[24.7.3.3] \textbf{Class template extents.} Specifies the facility family Class template extents. Read its declarations together with mandates, preconditions, effects, results, exceptions, complexity, and lifetime rules.
\item[24.7.3.3.1] \textbf{Overview.} Frames the terminology, applicability, and relationships needed to interpret Overview; detailed rules remain in the following subclauses.
\item[24.7.3.3.2] \textbf{Exposition-only helpers.} Specifies the facility family Exposition-only helpers. Read its declarations together with mandates, preconditions, effects, results, exceptions, complexity, and lifetime rules.
\item[24.7.3.3.3] \textbf{Constructors.} Specifies how Constructors establishes object state, including participation constraints, initialization effects, and exceptional exits.
\item[24.7.3.3.4] \textbf{Observers of the multidimensional index space.} Specifies the facility family Observers of the multidimensional index space. Read its declarations together with mandates, preconditions, effects, results, exceptions, complexity, and lifetime rules.
\item[24.7.3.3.5] \textbf{Comparison operators.} Specifies the facility family Comparison operators. Read its declarations together with mandates, preconditions, effects, results, exceptions, complexity, and lifetime rules.
\item[24.7.3.3.6] \textbf{Alias template dextents.} Specifies the facility family Alias template dextents. Read its declarations together with mandates, preconditions, effects, results, exceptions, complexity, and lifetime rules.
\item[24.7.3.4] \textbf{Layout mapping.} Specifies the facility family Layout mapping. Read its declarations together with mandates, preconditions, effects, results, exceptions, complexity, and lifetime rules.
\item[24.7.3.4.1] \textbf{General.} Frames the terminology, applicability, and relationships needed to interpret General; detailed rules remain in the following subclauses.
\item[24.7.3.4.2] \textbf{Requirements.} States the admissibility and semantic obligations for Requirements. Separate compile-time participation rules from caller preconditions and runtime effects.
\item[24.7.3.4.3] \textbf{Layout mapping policy requirements.} States the admissibility and semantic obligations for Layout mapping policy requirements. Separate compile-time participation rules from caller preconditions and runtime effects.
\item[24.7.3.4.4] \textbf{Layout mapping policies.} Specifies the facility family Layout mapping policies. Read its declarations together with mandates, preconditions, effects, results, exceptions, complexity, and lifetime rules.
\item[24.7.3.4.5] \textbf{Class template layout\_\allowbreak{}left::mapping.} Specifies the facility family Class template layout\_\allowbreak{}left::mapping. Read its declarations together with mandates, preconditions, effects, results, exceptions, complexity, and lifetime rules.
\item[24.7.3.4.5.1] \textbf{Overview.} Frames the terminology, applicability, and relationships needed to interpret Overview; detailed rules remain in the following subclauses.
\item[24.7.3.4.5.2] \textbf{Constructors.} Specifies how Constructors establishes object state, including participation constraints, initialization effects, and exceptional exits.
\item[24.7.3.4.5.3] \textbf{Observers.} Specifies the facility family Observers. Read its declarations together with mandates, preconditions, effects, results, exceptions, complexity, and lifetime rules.
\item[24.7.3.4.6] \textbf{Class template layout\_\allowbreak{}right::mapping.} Specifies the facility family Class template layout\_\allowbreak{}right::mapping. Read its declarations together with mandates, preconditions, effects, results, exceptions, complexity, and lifetime rules.
\item[24.7.3.4.6.1] \textbf{Overview.} Frames the terminology, applicability, and relationships needed to interpret Overview; detailed rules remain in the following subclauses.
\item[24.7.3.4.6.2] \textbf{Constructors.} Specifies how Constructors establishes object state, including participation constraints, initialization effects, and exceptional exits.
\item[24.7.3.4.6.3] \textbf{Observers.} Specifies the facility family Observers. Read its declarations together with mandates, preconditions, effects, results, exceptions, complexity, and lifetime rules.
\item[24.7.3.4.7] \textbf{Class template layout\_\allowbreak{}stride::mapping.} Specifies the facility family Class template layout\_\allowbreak{}stride::mapping. Read its declarations together with mandates, preconditions, effects, results, exceptions, complexity, and lifetime rules.
\item[24.7.3.4.7.1] \textbf{Overview.} Frames the terminology, applicability, and relationships needed to interpret Overview; detailed rules remain in the following subclauses.
\item[24.7.3.4.7.2] \textbf{Exposition-only helpers.} Specifies the facility family Exposition-only helpers. Read its declarations together with mandates, preconditions, effects, results, exceptions, complexity, and lifetime rules.
\item[24.7.3.4.7.3] \textbf{Constructors.} Specifies how Constructors establishes object state, including participation constraints, initialization effects, and exceptional exits.
\item[24.7.3.4.7.4] \textbf{Observers.} Specifies the facility family Observers. Read its declarations together with mandates, preconditions, effects, results, exceptions, complexity, and lifetime rules.
\item[24.7.3.5] \textbf{Accessor policy.} Specifies the facility family Accessor policy. Read its declarations together with mandates, preconditions, effects, results, exceptions, complexity, and lifetime rules.
\item[24.7.3.5.1] \textbf{General.} Frames the terminology, applicability, and relationships needed to interpret General; detailed rules remain in the following subclauses.
\item[24.7.3.5.2] \textbf{Requirements.} States the admissibility and semantic obligations for Requirements. Separate compile-time participation rules from caller preconditions and runtime effects.
\item[24.7.3.5.3] \textbf{Class template default\_\allowbreak{}accessor.} Specifies the facility family Class template default\_\allowbreak{}accessor. Read its declarations together with mandates, preconditions, effects, results, exceptions, complexity, and lifetime rules.
\item[24.7.3.5.3.1] \textbf{Overview.} Frames the terminology, applicability, and relationships needed to interpret Overview; detailed rules remain in the following subclauses.
\item[24.7.3.5.3.2] \textbf{Members.} Specifies the facility family Members. Read its declarations together with mandates, preconditions, effects, results, exceptions, complexity, and lifetime rules.
\item[24.7.3.6] \textbf{Class template mdspan.} Specifies the facility family Class template mdspan. Read its declarations together with mandates, preconditions, effects, results, exceptions, complexity, and lifetime rules.
\item[24.7.3.6.1] \textbf{Overview.} Frames the terminology, applicability, and relationships needed to interpret Overview; detailed rules remain in the following subclauses.
\item[24.7.3.6.2] \textbf{Constructors.} Specifies how Constructors establishes object state, including participation constraints, initialization effects, and exceptional exits.
\item[24.7.3.6.3] \textbf{Members.} Specifies the facility family Members. Read its declarations together with mandates, preconditions, effects, results, exceptions, complexity, and lifetime rules.
\end{description}
\section{Language-Lawyer and Library-Usage Pitfalls}
\begin{CornellReferencePitfallBox}\begin{itemize}[label=\textcolor{CornellReferenceAmber}{$\blacktriangleright$}]
\item Reading a synopsis as the complete behavioral contract.
\item Ignoring mandates, preconditions, exception behavior, or complexity requirements.
\item Using an exposition-only name as though it were public API.
\item Assuming invalidation, ownership, or thread-safety properties that are not stated.
\item Treating successful compilation as proof that semantic requirements and preconditions are met.
\end{itemize}\end{CornellReferencePitfallBox}
\section{Programmer and Implementation Checklist}
\begin{CornellReferenceChecklistBox}\begin{itemize}[label=$\square$]
\item Identify the declaring header or module exposure for each facility.
\item Check template arguments and mandates before evaluating an operation.
\item Establish every precondition at the call site.
\item Record effects, returned value, postconditions, and exceptions separately.
\item Audit iterator, reference, pointer, and object lifetime after mutation.
\item Verify stated complexity and synchronization assumptions.
\item For \CornellClauseRef{24.1}, verify general against the precise rule category and all cited dependencies.
\item For \CornellClauseRef{24.2}, verify requirements against the precise rule category and all cited dependencies.
\item For \CornellClauseRef{24.3}, verify sequence containers against the precise rule category and all cited dependencies.
\item For \CornellClauseRef{24.4}, verify associative containers against the precise rule category and all cited dependencies.
\end{itemize}\end{CornellReferenceChecklistBox}
\section{Clause Synthesis}
\begin{CornellReferenceOrientationBox}Defines sequence and associative containers, unordered containers, views and adaptors, container requirements, iterator invalidation, allocator use, complexity, and specialized containers such as arrays, spans, and multidimensional views. The standard model for owning and non-owning collections. A sound review distinguishes syntax and participation from runtime preconditions, keeps notes and examples non-mandatory, and follows cross-clause dependencies before making a portability claim.\end{CornellReferenceOrientationBox}
\section{Self-Test Questions}
\begin{enumerate}
\item What role does General play in the clause's overall model?
\item Which normative distinctions must be kept separate when applying Requirements?
\item How would you audit a program or implementation that depends on Sequence containers?
\item Compare Associative containers with Unordered associative containers; what different question does each answer?
\item What role does Unordered associative containers play in the clause's overall model?
\item Which normative distinctions must be kept separate when applying Container adaptors?
\item How would you audit a program or implementation that depends on Views?
\item Compare Preamble with General containers; what different question does each answer?
\item What role does General containers play in the clause's overall model?
\item Which normative distinctions must be kept separate when applying Container data races?
\item How would you audit a program or implementation that depends on Sequence containers?
\item Compare Node handles with Insert return type; what different question does each answer?
\item What role does Insert return type play in the clause's overall model?
\item Which normative distinctions must be kept separate when applying Associative containers?
\item Scenario: a program relies on an unstated implementation behavior while using Unordered associative containers. What must be checked before calling the program portable?
\item Scenario: a program relies on an unstated implementation behavior while using General. What must be checked before calling the program portable?
\end{enumerate}
\clearpage\section{Answer Key}
\begin{enumerate}
\item General is one part of the clause's library semantics model. Its role is understood by connecting the local rule in 24.7.3.5.1 to the clause orientation and its dependent concepts.
\item Keep formation and well-formedness separate from runtime preconditions and effects. Also distinguish mandatory wording from notes, implementation choice from undefined behavior, and interface declarations from detailed contracts; see 24.7.3.5.2.
\item Audit Sequence containers by locating the controlling wording in 24.3, listing inputs and type constraints, proving any preconditions, recording effects and failure behavior, and checking lifetime, invalidation, complexity, and synchronization where applicable.
\item Associative containers and Unordered associative containers occupy different steps or facility groups. Analyze each under its own subclause, then follow the explicit dependency between them rather than transferring one set of guarantees to the other; begin at 24.4.
\item Unordered associative containers is one part of the clause's library semantics model. Its role is understood by connecting the local rule in 24.5 to the clause orientation and its dependent concepts.
\item Keep formation and well-formedness separate from runtime preconditions and effects. Also distinguish mandatory wording from notes, implementation choice from undefined behavior, and interface declarations from detailed contracts; see 24.6.
\item Audit Views by locating the controlling wording in 24.7, listing inputs and type constraints, proving any preconditions, recording effects and failure behavior, and checking lifetime, invalidation, complexity, and synchronization where applicable.
\item Preamble and General containers occupy different steps or facility groups. Analyze each under its own subclause, then follow the explicit dependency between them rather than transferring one set of guarantees to the other; begin at 24.2.1.
\item General containers is one part of the clause's library semantics model. Its role is understood by connecting the local rule in 24.2.2 to the clause orientation and its dependent concepts.
\item Keep formation and well-formedness separate from runtime preconditions and effects. Also distinguish mandatory wording from notes, implementation choice from undefined behavior, and interface declarations from detailed contracts; see 24.2.3.
\item Audit Sequence containers by locating the controlling wording in 24.3, listing inputs and type constraints, proving any preconditions, recording effects and failure behavior, and checking lifetime, invalidation, complexity, and synchronization where applicable.
\item Node handles and Insert return type occupy different steps or facility groups. Analyze each under its own subclause, then follow the explicit dependency between them rather than transferring one set of guarantees to the other; begin at 24.2.5.
\item Insert return type is one part of the clause's library semantics model. Its role is understood by connecting the local rule in 24.2.6 to the clause orientation and its dependent concepts.
\item Keep formation and well-formedness separate from runtime preconditions and effects. Also distinguish mandatory wording from notes, implementation choice from undefined behavior, and interface declarations from detailed contracts; see 24.4.
\item First identify the exact requirement in 24.5, then classify the relied-on behavior as required, unspecified, implementation-defined, conditionally supported, or undefined. Portability requires satisfying the program obligations and avoiding dependence on a choice not guaranteed in the target environments.
\item First identify the exact requirement in 24.7.3.5.1, then classify the relied-on behavior as required, unspecified, implementation-defined, conditionally supported, or undefined. Portability requires satisfying the program obligations and avoiding dependence on a choice not guaranteed in the target environments.
\end{enumerate}
\section{Key-Term Recap}
\begin{longtable}{@{}>{\RaggedRight\arraybackslash}p{0.28\textwidth}>{\RaggedRight\arraybackslash}p{0.64\textwidth}@{}}\toprule\textbf{Term} & \textbf{Working meaning}\\\midrule\endfirsthead\toprule\textbf{Term} & \textbf{Working meaning (continued)}\\\midrule\endhead\bottomrule\endfoot
container & A type that stores or exposes a collection through common requirements.\\
iterator invalidation & Loss of validity of iterators, pointers, or references after an operation.\\
allocator-aware container & A container whose storage and construction follow allocator protocols.\\
contiguous container & A container whose elements occupy one contiguous memory region.\\
node handle & A movable ownership object for an extracted associative-container node.\\
span & A non-owning contiguous sequence view.\\
mdspan & A multidimensional view combining data handle, mapping, and accessor.\\
\end{longtable}
\CornellTakeaway{Container use is correct only when requirements, complexity, allocator behavior, and invalidation are all tracked per operation.}
\end{document}
